\section{结构与协议证明的补充细节}\label{app:protocols}
\subsection{原子来源组的附加参数}\label{app:atomic}
第~\ref{sec:frontier}节的二元候选边界不能直接覆盖原子 RF 超边。若另有 $d=\max_r|D(r)|$，把每个原子组 $\{r\}\cup D(r)$ 安排在其最晚端点引入时处理，所有端点连续保留到该位置。考虑切口处一个因该超边而额外存活、却没有二元未来邻居的旧写 $w$。组中心读 $r$ 必在过去，否则 $w$ 就有未来邻居 $r$；该组又必有某个未来候选写 $v$。由 R1 饱和，$v\in D(r)$ 意味着 $\neg H(r,v)$，所以 $r\in U_t$。每个这样的读最多贡献 $d$ 个额外旧写。因此在定理~\ref{thm:frontier} 的袋上增加至多 $cd$ 个端点，即得
\[
 K_{\rm atomic}\le c(d+1)+ap+p+a+1.
\]
按最后端点调度且一直保留端点保证超边覆盖和连续生存期。这个结论以 $d$ 为附加参数；若它随历史增长，则不能得到仅依赖 $c,p,a$ 的统一常数界。

\subsection{无界纯失败读链的完整袋展开}\label{app:chains}
对第~\ref{sec:buffer}节程序，识别由已观测失败分支确定的极大连续失败尝试串。串内每次尝试只有两个索引读，没有载荷事件、写或其他地址访问。引理~\ref{lem:bufferwindow} 及自身索引唯一性给出阶段 $k$ 的固定来源锚：生产者失败串使用 $W_{k-1}$ 与 $A_{k-1-C}$，消费者失败串使用 $W_{k-1}$ 与 $A_{k-1}$。负下标对应的失败分支已被证明不可能；初始化锚另计。

规划骨架保留全部成功操作和每线程最后尚未决定分支的调用，每骨架阶段最多 $b_0=2(B+3)$ 个事件。骨架只用于安排袋；所有 $H/L$ 和来源域仍从完整原始输入证明得到。暂时缩去失败链，不等于删除输入事件后重新解释一个协议。特别是被缩去的链首尾所形成的规划连接不作为额外 $\SB$ 语义边加入后端。

为每串登记来源锚及左右真实线程端点，在其所属阶段 $k$ 设置处理位置。把已引入的阶段 $[k-C-1,k]$ 骨架事件连续保留到该位置及阶段末尾。该窗口有 $C+2$ 个阶段，所有非初始化锚和左右端点都位于其中；左端点属于阶段 $k-1$ 或 $k$，右端点若存在则属于 $k$。若完整输入以纯失败串结束而该阶段没有骨架事件，则增加一个仅用于计划的虚拟处理位置，缺失右端点不补造。

对于窗口之前仍有实际候选连接的骨架事件，采用定理~\ref{thm:frontier} 的逐边论证。阶段分隔使旧读不再具有未来 RF，旧写与未来写的未决 CO 不存在；若旧稳定写不为稳定极大写，R2 删除其未来 RF，且中间稳定写阻止它成为 CO 覆盖前驱。它们至多贡献 $ap$ 个写。骨架线程连接另需至多 $p$ 个尾端点，初始化需 $a$ 个。实际候选边界及登记位置决定生存期；稳定极大集合只用于计数，不能让已经退休事件因后来成为极大元素而重新出现。因此所有锚同时可用时，骨架袋大小至多
\[
 (C+2)b_0+ap+p+a.
\]

在该袋处一次展开一条串：引入第一个真实读，加入其固定 RF 及左侧真实线程边；逐个引入下一读，处理它的固定 RF 和前一读到它的线程边，再退休前一读；最后处理末读到右端点的边。最多同时增加两个读。原袋端点在整个展开期间连续保留，新读也只在连续区间出现；不同串顺次展开。于是 $K\le(C+2)b_0+ap+p+a+2$，代入 $p=2$ 得式~\eqref{eq:chainbound}。

骨架袋允许右线程端点先于中间读引入。这是图处理顺序，真正加入语义的仍只有每条原始邻接边与所有真实读。读的固定正 RF 以及所有被排除候选的负 RF 保留唯一标签归属；监视器仍经过这些读计算 HB、FR 和禁止游走。因此展开保持完整投影，而不是只保持删除失败读后的一个商输入。若重试体还含其他工作，这个纯链证明不适用，应在阶段工作量参数中计入那些事件。

\subsection{屏障失败等待和未完成前缀}\label{app:barriertails}
引理~\ref{lem:barrier} 的成功对齐还确定到达地址的 CO。$A_{t,i}\mathrel\rf R_{t,i}\mathrel\SB Z_{t,i}$ 给出 $A_{t,i}<_{\co}Z_{t,i}$；$Z_{t,i}\mathrel\SB G_{t,i}\mathrel\rf S_{t,i}\mathrel\SB A_{t,i+1}$ 给出 $Z_{t,i}<_{\co}A_{t,i+1}$。初始化最小，形成正文所列交错写链，截断前缀只会删除末尾而不会插入未观测事件。

证明正文布局中失败读不属于前沿，需要分别检查两类等待。协调者调用 $i$ 中对 $a_t$ 的失败读 $f$，若本轮到达写 $A_{t,i}$ 存在，已经在布局的工作者步骤引入。本轮及以后重置 $Z$ 在线程中都位于 $f$ 之后。未来工作者到达 $A_{t,j}$、$j>i$ 若存在，则其前缀和成功对齐要求 $S_{t,i}$ 存在，并给出
\[
 f\mathrel{\SB^+}G_{t,i}\mathrel\rf S_{t,i}\mathrel{\SB^+}A_{t,j}.
\]
若当前等待未成功且所需到达事件缺失，合法前缀就不能含需要它的后续唤醒及更晚调用。这不是根据调度假定其发生，而是由非空成功来源域导出的必要存在性。因此所有仍在未来的同地址写均在 $H(f,\cdot)$ 后方。

工作者调用 $i$ 的失败唤醒读 $f$ 访问 $g_t$。本轮 $G_{t,i}$ 若存在，已在协调者步骤引入；未来只能有 $G_{t,j}$、$j>i$。其存在要求 $R_{t,j}$ 及准确来源 $A_{t,j}$ 存在，因此
\[
 f\mathrel{\SB^+}A_{t,j}\mathrel\rf R_{t,j}\mathrel{\SB^+}G_{t,j}.
\]
这说明 $f$ 对所有未来同地址写稳定。若不存在未来写，稳定条件自然成立。读与未来同地址读不构成式~\eqref{eq:frontier} 中的相关对。于是全部失败读均被排除出 $U_t$，它们的线程边仍由通用界的尾事件项处理。

任意阶段 $i$ 的工作者事件在其下一到达 $A_{t,i+1}$ 之前；下一到达通过成功读取进入协调者，再经其全部到达之后的唤醒进入任意线程的阶段 $i+2$。协调者阶段 $i$ 的事件直接沿程序顺序到下一轮唤醒。只要目标阶段 $i+2$ 的事件存在，所有这些必要交接都由前缀和成功对齐强制存在。因此阶段差至少二的事件均有共同 $H$。去掉已证明稳定的失败读，切口最多剩两个阶段各 $b_0$ 个骨架事件，得到式~\eqref{eq:barrierbound}。

\subsection{地址封闭组合的证明}\label{app:composition}
设地址分区为 $\mathcal A_1,\ldots,\mathcal A_h$，$E_i$ 包含访问 $\mathcal A_i$ 的全部事件及相应初始化。局部线程顺序是全局线程序在 $E_i$ 上的限制；它的邻接边可以跨过其他组件的事件。局部协议元数据必须与这个控制投影兼容；观测约束采用由全局输入蕴含的局部限制。

任取全局合法完成 $J$。局部读的真实来源因为同地址闭合仍在 $E_i$ 内，局部 CO 是原地址的完整写序；局部线程边对应全局线程路径，所以局部 $\hb$ 包含于全局 $\hb$。局部 $\eco$ 也包含于全局 $\eco$。若局部出现禁止闭环，全局也出现同一闭环，矛盾。因此 $J|E_i$ 合法，局部已证 $H_i$ 在全局中可靠。将全部 $H_i$ 与全局 $\SB$、初始化取传递闭包，仍为所有全局完成的共同下界。

建立布局约束图 $T$，加入每个局部总排列 $\theta_i$ 的邻接边及全局线程邻接。若 $T$ 无环，固定一个拓扑序 $\theta$。由于每个 $\theta_i$ 为总序，$\theta$ 在 $E_i$ 上的投影恰为 $\theta_i$；全局每个切口对应局部的某个切口。取 $e\in U_t$ 及其未来相关见证 $f$，二者因同地址而属于同一个 $E_i$。若 $e\notin U_i$，局部稳定性给出 $H_i(e,f)$，进而给出全局 $H(e,f)$，与 $e\in U_t$ 矛盾。所以 $U_t\subseteq\bigcup_i U_i$，得到 $c\le\sum_i c_i$ 和正文组合界。

地址闭合不可省略。若把两个各自只写同一个 $x$ 的线程分成两个“组件”，每个局部前沿均为零；全局却有未决 CO 的完全二部候选，前沿可随写数增长。相容性也不可省略：局部总排列与跨组件线程边可能有环，表示给定计划无法合并，不能由此判断 RA 观测无解。对于无界失败缓冲，式~\eqref{eq:chainbound} 只给袋大小而非原始前沿界，故不能把该 $K$ 当作局部 $c_i$ 代入组合定理。可直接组合的是实际已给 $\SB$ 相容布局和前沿界的局部组件。

\subsection{完整交替发布网络的无写覆盖推论}\label{app:publication}
为明确第~\ref{sec:frontier}节的推论范围，设 $p$ 个线程构成直径为 $d$ 的固定强连通有向网络。每线程独占一个初值 $0$ 的标志；共有同样的 $h_0$ 个完整阶段，阶段 $i$ 向本线程标志写至少两次相同值 $i\bmod2$，再对每条入边读对方标志，观测 $(i-1)\bmod2$。这些是标志的全部访问，标志与数据地址分离，每阶段全部线程事件数至多 $b$。

对一条网络边，接收线程的 $h_0$ 个读按程序顺序排列。单写者 CO 和 RR 一致性使来源块号非降；交替读值使块号严格增加。最后一个读的值与最后一个写块不同，因此其块号至多 $h_0-1$。$h_0$ 个非负严格递增块号只能为 $0,1,\ldots,h_0-1$，其中 $0$ 是初始化。故第 $i$ 次等待的来源锁在对方第 $i-1$ 块，但块内身份仍可歧义。第一次发布是该块全部候选的共同前序，使用 $R_\forall$ 即得共同交接。

线程 $t$ 阶段 $i$ 的任意事件先于它阶段 $i+1$ 的首次发布，该发布又先于邻居阶段 $i+2$ 的对应等待。沿长度至多 $d$ 的网络路径重复这一步，再到目的线程下一阶段，可得阶段差至少 $2d+1$ 的共同因果分隔。因此 $c\le(2d+1)b$，无需规定数据地址的写频率。

这一改进有具体的参数分离族。在正文两线程发布内核中增加地址 $z$：A 只在第一阶段写 $z=1$，中间每阶段读 $z=1$，最后阶段再写 $z=1$。$p,a,b,d$ 均可固定，而最长无写区间随阶段数增加。两个非初始化写形成一条跨越整个历史的 CO 覆盖边，所有中间读仍来自首写；它们本就在同一线程顺序中，因而不为旧事件增加跨线程 HB，且自身一致性直接成立。一个稳定旧写可以承载这些长边。这里分离的是统一参数界，不声称每个有限输入都不能选取某个随规模增长的有限写覆盖参数。

\subsection{同值发布例子的完整见证}\label{app:witness}
为第~\ref{sec:example}节补齐省略事件。$x,y,fa,fb$ 初始化均为 $0$。对 $i=1,2$，两个线程依次执行
\begin{align*}
 A_i:\ &A_i{:}x=\W(x,1);\ A_i{:}0=\W(fa,i\bmod2);\ A_i{:}y=\W(y,1);\\
        &A_i{:}1=\W(fa,i\bmod2);\ A_i{:}w=\R(fb,(i-1)\bmod2);\ A_i{:}r=\R(x,1),\\
 B_i:\ &B_i{:}y=\W(y,1);\ B_i{:}0=\W(fb,i\bmod2);\ B_i{:}x=\W(x,1);\\
        &B_i{:}1=\W(fb,i\bmod2);\ B_i{:}w=\R(fa,(i-1)\bmod2);\ B_i{:}r=\R(y,1).
\end{align*}
正文的 $a_0,u,a_1,v,r,s$ 分别为 $A_1{:}0,A_1{:}y,A_1{:}1,B_2{:}y,B_2{:}w,B_2{:}r$。所有数据读选择本线程本轮的首个数据写，第一轮两个等待均选择初始化。反向 CO 见证取第二轮等待都读对方第一轮早发布，并取
\begin{align*}
 x:&\quad i_x<B_1{:}x<B_2{:}x<A_1{:}x<A_2{:}x,\\
 y:&\quad i_y<B_1{:}y<B_2{:}y<A_1{:}y<A_2{:}y.
\end{align*}
标志 CO 为初始化之后按各自线程顺序。它满足 $d\land e\land\neg q$。所有跨线程非初始化 RF 只到第二轮的末尾等待；因此 HB 无环。对于数据地址，来自对方早发布之前的数据写只能到达另一个地址的数据读：$A_1{:}x$ 到 $B_2{:}r$，$B_1{:}y$ 到 $A_2{:}r$。不存在与反向 CO 冲突的同地址跨线程 HB。各线程内部的同地址来源与写序依次前进，满足 WW／WR／RW／RR。标志读选择的早发布之前没有更晚同地址 HB 前驱写，其未来同地址写在后续块；同样满足四类一致性。因此完整赋值合法。

正向见证把 $x,y$ 的 CO 都按轮拼接，轮内两写任意选择；第二轮等待可选择对方第一轮的早或晚发布。所有跨线程 HB 严格增轮，同轮仅有线程内部 HB，来源为本线程本轮写。由此 WW 与轮序相容，WR 的同轮同地址前驱至多为实际来源，RW 的后继写在以后轮，RR 来源轮号非降；标志按块顺序一致。分别选择 B 的早或晚来源，得到 $d\land e\land q$ 及 $d\land l\land q$ 的完整合法完成。正文对 $d\land l\land\neg q$ 的 WR 反证排除了最后一种组合。
